\documentclass[11pt,a4paper]{article}
\usepackage[margin=23mm]{geometry}
\usepackage{amsmath,amssymb,amsthm}
\setlength{\emergencystretch}{2em}
\newtheorem{lemma}{Lemma}
\newtheorem{theorem}[lemma]{Theorem}
\newcommand{\gen}[1]{\langle#1\rangle}
\newcommand{\Sym}{\operatorname{Sym}}
\title{The size of an orbit with a local crossing}
\author{Human Compression project: L044 continuation}
\date{}
\begin{document}\maketitle
All marked generators below are involutions.
\begin{lemma}[An anchored action]
Let $T=\gen{t_0,\ldots,t_{h-1}}$ be transitive on $V$,
where the $t_i$ are involutions with string commutations,
$h\ge1$, and $t_{h-1}\ne1$.
Suppose $a\in V$ is fixed by $t_1,\ldots,t_{h-1}$.
Then $|V|\ge h+1$. In equality $T=\Sym(V)$.
\end{lemma}
\begin{proof}
For $h=1$ this is the two-point orbit of a nontrivial
involution. For $h\ge2$, choose a shortest path from
$a$, using $t_0,\ldots,t_{h-2}$, to a point moved
by $t_{h-1}$. Such a point exists, or the prefix
orbit would be invariant under $T$ and the last
mark would fix all of $V$. Every preceding point
on this path is fixed by the last mark, so its
final edge adds a new point. The complete acquired
no-gap argument makes every label $0,\ldots,h-2$
occur: at a missing index, commute the higher
factors, which fix $a$, past the lower factors,
which commute with the last mark. Thus the path,
including its final $t_{h-1}$-edge, has at least
$h$ edges, proving the bound.
In equality its labels occur exactly once in
increasing order. The complete acquired exhausted
anchored-path lemma gives the natural symmetric
group when $h\ge3$; for $h=2$ the two distinct
transpositions on three points generate $S_3$.
No independence or IP is required.
\end{proof}

\begin{lemma}[Two commuting anchored factors]
Suppose $L,R$ commute and $LR$ is transitive on $A$.
At $a\in A$, the nonidentity left restrictions have
string indices ending at $j-1$, all except that
last mark fixing $a$; the right restrictions start
at $j+1$, all except that first mark fixing $a$.
Let $h,u$ be their respective numbers of active
labels. Then
\[
 |A|\ge h+u+1,
\]
except possibly when $|A|=2$, $h=u=1$, and both
factors induce the same transposition on $A$.
An empty factor is permitted and has zero labels.
\end{lemma}
\begin{proof}
If $A$ is a singleton, both factors are trivial
and the bound holds. Hence assume $|A|\ge2$.
The active labels within each factor form one
interval adjoining $j$. Indeed any other marked
component fixes $a$ and commutes with the remaining
components, so is trivial on the transitive orbit
of its own factor. Commutation with the opposite
factor then makes it trivial on all of $A$.
Any nonidentity restriction in a factor is also
nonidentity on its orbit through $a$: if it fixed
that orbit pointwise, commuting with the opposite
factor would make it trivial on $A$.
The preceding lemma, in reversed order on the left,
gives $|a^L|\ge h+1$ and $|a^R|\ge u+1$ for
nonempty factors.

If neither factor is transitive, its orbits have
equal sizes: the opposite commuting factor
permutes them transitively. Hence
\[
 |A|\ge2\max\{|a^L|,|a^R|\}\ge h+u+2.
\]
If $L$ is transitive, every right mark other than
the one of index $j+1$ centralizes $L$ and fixes
$a$, so is trivial. Thus $u\le1$. For $u=0$ the
anchored bound suffices. For $u=1$ the centralizer
of $L$ on $A$ is nontrivial. If $h\ge2$, equality
$|A|=h+1$ would make $L$ natural symmetric of
degree at least three, whose centralizer is
trivial. Therefore $|A|\ge h+2=h+u+1$.
If $h=1$, transitivity of its single involution
forces $|A|=2$, giving precisely the stated
exception. The case with $R$ transitive is the
same with the factors interchanged.
\end{proof}

\begin{theorem}[The whole local orbit]
Let $H=\gen{\rho_m,\ldots,\rho_x}$ act transitively
on $O$, with string commutations and every displayed
restriction nonidentity. Suppose the full local
omission of some $j\in[m,x]$ has exactly two
orbits $A,B$, joined by a unique $j$-edge $\{a,b\}$.
Assume also that the omission of the first mark
$\rho_m$ is intransitive on $O$. Then
\[
                            |O|\ge x-m+2.
\]
\end{theorem}
\begin{proof}
Put $L=H_{<j}$ and $R=H_{>j}$. They commute and
$LR$ is transitive on each of $A,B$. Every distant
mark fixes $a,b$: it preserves the local omission
orbits and the unique crossing edge. Thus the
two preceding anchored factors apply on each side.
Write their active lengths as $h_A,u_A$ and
$h_B,u_B$, taking a zero for an empty factor.
The active intervals adjoining $j$ proved above
give
\[
      \max(h_A,h_B)=j-m=:h,\qquad
      \max(u_A,u_B)=x-j=:u.
\]
If these maxima are attained on different sides,
the individual anchored bounds give
$|A|+|B|\ge(h+1)+(u+1)$.
If they are both attained on one side, say $A$,
the preceding lemma gives $|A|\ge h+u+1$ and
the nonempty other side adds one more point.

The only unaccounted possibility is $h=u=1$,
$|A|=2$ and $|B|=1$. Then the first and last
marks are the same transposition on $A$, both
fixing $B$. The middle mark is the crossing
transposition on the three points. Omitting the
first mark leaves a transitive action, contrary
to the hypothesis. If $|B|\ge2$, the desired
bound already holds. This proves all cases,
including an empty left or right factor.
\end{proof}

\begin{theorem}[When the central involution exchanges the witness orbit]
Retain every actual odd-degree head hypothesis of
\emph{Local pairs and ambient witnesses}, including
two ambient eligible edges globally for every
tail colour. Let $O$ be its normalized orbit with
a genuine local unique crossing, with active
interval $\{5,\ldots,x\}$. If $\delta$ exchanges
$O$ with a distinct $H$-orbit, then
\[
                x<r-1,\qquad
                |\operatorname{Fix}(G_{>x})|\ge2x+1.
\]
\end{theorem}
\begin{proof}
The first omission on $O$ is $K=G_{>5}$ and is
intransitive by normalization. The preceding
theorem gives $|O|\ge x-3$. Since $\delta$
centralizes $H$, its other orbit $O'$ has the
same complete marked action and the same size.
All marks after $x$ fix both orbits. They also
fix the seven points $D\cup\{b\}$, disjoint
from them. Therefore
\[
            |\operatorname{Fix}(G_{>x})|
                         \ge2(x-3)+7=2x+1.
\]

Suppose $x=r-1$. The acquired actual anchor
$z=b^{\rho_4}$ is fixed by $\delta$ and by $K$.
It is outside $D,b$ and is not $H$-fixed: otherwise
the full group would preserve $D\cup\{b,z\}$,
as in the complete last-fixed-point proof.
It therefore belongs to a nontrivial $H$-orbit
$V$ disjoint from $O,O'$ and the seven points.
The equation $n=2r+1$ and $|O|=|O'|\ge r-4$
leave at most two points for $V$. Consequently
$|O|=r-4$, $|V|=2$, and these sets exhaust $\Omega$.
The group $K$ fixes $V$ pointwise, since it fixes
its point $z$. Transitivity of $H|_V$ forces
$\rho_5$ to exchange the two points of $V$.
Its restrictions on $O,O'$ are conjugate by
$\delta$ and therefore have the same sign.
It fixes $D,b$. Thus $\rho_5$ is odd on all
of $\Omega$, a contradiction. The terminal
case is excluded.
\end{proof}

\begin{theorem}[The witness orbit is fixed pointwise]
Under the preceding actual head hypotheses, $\delta$
fixes $O$ pointwise.
\end{theorem}
\begin{proof}
First suppose $\delta$ exchanges $O,O'$. The preceding
proof supplies $x<r-1$ and the point $z=b^{\rho_4}$,
which is fixed by $\delta,K$ but not by $H$.
Thus $y=z^{\rho_5}\ne z$. Both lie in the same
$H$-orbit, disjoint from $O,O'$ and $D\cup\{b\}$.

If $x\ge6$, every mark after $x$ commutes with
$\rho_5$ and fixes $z$, so fixes $y$ as well.
Together with $O,O'$ and the seven head points,
these give at least $2(x-3)+9=2x+3$ fixed points.
Apply the complete ambient interval bound to
$I=\{x+1,\ldots,r-1\}$. It is nonempty, every
colour has two ambient eligible edges, and its
action is intransitive because of these fixed
points. If $U$ is its support and $c\ge1$ its
number of nontrivial orbits, that bound gives
\[
 |U|\ge2(r-1-x)+c\ge2(r-1-x)+1.
\]
The disjoint fixed points and support then force
$n\ge2r+2$, contrary to $n=2r+1$.

It remains that $x=5$. The transitive action on
$O$ is a single nontrivial involution, so
$|O|=|O'|=2$. Since $x<r-1$, the interval
$K=G_{>5}$ is nonempty. It fixes the twelve
points
\[
             D\cup\{b\}\cup O\cup O'\cup\{z\}.
\]
The same interval bound gives $f+c\le13$, where
$f$ counts its fixed points and $c\ge1$ its
nontrivial orbits. Therefore $f=12$, $c=1$;
its unique nontrivial orbit $U$ has odd size
$n-12=2r-11$. The involution $\delta$ centralizes
$K$, so preserves $U$. A nonidentity centralizer
element of a transitive action fixes no point;
an involution with no fixed point has even
support. Since $|U|$ is odd, $\delta|_U=1$.
On the twelve fixed points it fixes $D,b,z$ and
exchanges $O$ with $O'$, acting as exactly two
transpositions. It is therefore even everywhere,
contrary to its stated oddness. This excludes
all orbit exchanges.

Hence $\delta$ preserves $O$. The actual
$\tau\in G_j$ induces $\delta$ there and
centralizes its full $H$-action. It preserves
the unique ambient witness edge, and ambient
eligibility prevents it from exchanging that
edge's endpoints. It fixes a point of a
transitive $H$-orbit and centralizes $H$, so
fixes the entire orbit pointwise.
\end{proof}

\paragraph{Attribution and remaining route.}
The anchored no-gap and exhausted-path proofs and
full ambient interval bound are acquired. The two
commuting-factor size count, terminal parity argument
and twelve-fixed-point parity argument are the new
steps. They avoid asserting disjointness of two point
paths. The actual witness orbit and its normalized
first omission come from the complete preceding proof.
The remaining route is now the orbit fixed pointwise
by $\delta$, together with full source routing and
global attachment. No new induction, computation,
formal, original or question credit is claimed.
\end{document}
